\idxletter{N}
\sourcepage{767}
\idxentry{0}{单例设计}{N-of-1 design}{226, 416g}
\idxentry{0}{全国门诊医疗调查（NAMCS）}{National Ambulatory Medical Care Survey (NAMCS)}{294}
\idxentry{0}{美国国家卫生统计中心（NCHS）}{National Center for Health Statistics (NCHS)}{294}
\idxentry{0}{全国健康与营养调查（NHANES）}{National Health and Nutrition Examination Survey (NHANES)}{31, 116–117, 293–294}
\idxentry{0}{美国国立卫生研究院（NIH）}{National Institutes of Health (NIH)}{21, 295}
\idxentry{0}{全国调查}{National surveys}{291, 293–294}
\idxentry{0}{自然实验}{Natural experiment}{168, 416g}
\idxentry{0}{NCHS。参见美国国家卫生统计中心}{NCHS. See National Center for Health Statistics (NCHS)}{}
\idxentry{0}{阴性预测值}{Negative predictive value}{247}
\idxentry{0}{新生儿血液标本研究}{Neonatal blood specimens, research on}{100–101}
\idxentry{0}{巢式病例队列设计}{Nested case-cohort design}{148–149}
\idxentry{1}{血浆ACE2与死亡或心脏代谢疾病风险}{plasma ACE2 and risk of death or cardiometabolic diseases}{148–149}
\idxentry{1}{优点}{strengths}{149}
\idxentry{1}{缺点}{weaknesses}{149}
\idxentry{0}{巢式病例对照研究}{Nested case-control study}{48, 138, 145, 146f, 246, 416g}
\idxentry{1}{优点}{strengths}{149}
\idxentry{1}{缺点}{weaknesses}{149}
\idxentry{0}{巢式双（或多）队列研究}{Nested double (or multiple) cohort study}{127, 128, 131, 416g}
\idxentry{0}{建立联系网络}{Networking}{284–285}
\idxentry{0}{问题的中立性}{Neutrality, of questions}{310}
\idxentry{0}{新研究者}{New investigators}{353}
\idxentry{0}{新技术的应用}{New technologies, application of}{18}
\idxentry{0}{NHANES。参见全国健康与营养调查}{NHANES. See National Health and Nutrition Examination Survey (NHANES)}{}
\idxentry{0}{NIH。参见美国国立卫生研究院}{NIH. See National Institutes of Health (NIH)}{}
\idxentry{0}{NNT。参见需治疗人数}{NNT. See Number needed to treat (NNT)}{}
\idxentry{0}{不追加经费延期}{No-cost extension}{326, 416g}
\idxentry{0}{无序分类变量}{Nominal variable}{40, 416g}
\idxentry{0}{非病例报告表数据}{Non-CRF data}{346}
\idxentry{0}{非差异性错分偏倚}{Nondifferential misclassification bias}{143, 416g}
\idxentry{0}{非劣效性}{Noninferiority}{221–224, 223f}
\idxentry{0}{非劣效界值}{Noninferiority margin}{223, 416g}
\idxentry{0}{非劣效试验}{Noninferiority trials}{74–75, 416g}
\idxentry{0}{非概率样本}{Nonprobability sample}{31, 33, 416–417g}
\idxentry{0}{非随机干预设计}{Nonrandom interventional designs}{229–230}
\idxentry{1}{非随机组内设计}{nonrandomized within-group designs}{232–233}
\idxentry{2}{前后试验设计}{before-after trial designs}{230–231, 230f}
\idxentry{2}{中断时间序列设计}{interrupted time series designs}{232, 232f}
\idxentry{0}{未应答偏倚}{Nonresponse bias}{34, 417g}
\idxentry{0}{未达统计显著的结果}{Nonsignificant result}{58}
\idxentry{0}{规范化}{Normalization}{342, 417g}
\sourcepage{768}
\idxentry{0}{研究问题的新颖性}{Novelty, research question and}{21}
\idxentry{0}{零假设}{Null hypothesis}{54, 55, 58, 64, 65, 417g}
\idxentry{0}{需伤害人数}{Number needed to harm}{417g}
\idxentry{0}{需治疗人数（NNT）}{Number needed to treat (NNT)}{83, 417g}
\idxentry{0}{数值变量}{Numerical variable}{40–41, 417g}
\idxentry{0}{养老机构住民研究}{Nursing home residents, research on}{104}
\idxletter{O}
\idxentry{0}{客观性}{Objectivity}{47}
\idxentry{0}{观察性研究}{Observational study}{3, 53, 116, 417g}
\idxentry{1}{优缺点}{advantages and disadvantages of}{151–152t}
\idxentry{1}{设计的选择}{choosing among}{150}
\idxentry{1}{检查对结局作用的研究}{studies of effect of testing on outcomes}{255–256}
\idxentry{0}{观察者偏倚}{Observer bias}{44, 417g}
\idxentry{0}{观察者变异}{Observer variability}{41, 417g}
\idxentry{0}{比值}{Odds}{131, 132t, 153, 155, 417g}
\idxentry{0}{比值比（OR）}{Odds ratio (OR)}{70, 132, 140, 417–418g}
\idxentry{1}{病例对照研究中的比值比}{for case-control studies}{153–154}
\idxentry{1}{与风险比的关系}{risk ratio and}{155}
\idxentry{0}{单样本$t$检验}{One-sample t test}{80, 418g}
\idxentry{0}{单侧假设}{One-sided hypothesis}{58, 418g}
\idxentry{0}{一对多}{One-to-many}{341, 418g}
\idxentry{0}{在线资料收集表}{Online data collection forms}{345}
\idxentry{0}{开放式问题}{Open-ended question}{47, 263, 307–308, 418g}
\idxentry{0}{操作手册}{Operations manual}{42, 50, 323, 327–328, 332, 418g}
\idxentry{1}{随机试验操作手册目录}{table of contents for randomized trial}{338}
\idxentry{0}{阿片过量死亡}{Opiate overdose deaths}{147–148}
\idxentry{0}{机会性研究设计}{Opportunistic study design}{165, 165t, 168, 178, 229, 418g}
\idxentry{1}{双重差分设计}{difference in differences designs}{170–171}
\idxentry{1}{工具变量}{instrumental variables}{169}
\idxentry{1}{中断时间序列设计}{interrupted time series designs}{170}
\idxentry{1}{孟德尔随机化}{Mendelian randomization}{168–169}
\idxentry{1}{自然实验}{natural experiments}{168}
\idxentry{1}{断点回归设计}{regression discontinuity designs}{169–170}
\idxentry{0}{乐观先验分布}{Optimistic prior distribution}{235, 418g}
\idxentry{0}{OR。参见比值比}{OR. See Odds ratio (OR)}{}
\idxentry{0}{有序分类变量}{Ordinal variable}{40, 418g}
\idxentry{0}{骨关节炎倡议}{Osteoarthritis Initiative}{296}
\idxentry{0}{男性骨质疏松性骨折研究}{Osteoporotic Fractures in Men Study}{296}
\idxentry{0}{结局}{Outcome}{418g}
\idxentry{1}{频率指标}{frequency, measures of}{131–132}
\idxentry{1}{关注的结局}{of interest}{4}
\sourcepage{769}
\idxentry{1}{测量}{measurements}{198–199}
\idxentry{2}{不良事件}{adverse events}{199–200}
\idxentry{2}{复合结局}{composite outcomes}{199}
\idxentry{2}{结局变量数量}{outcome variables, number of}{199}
\idxentry{1}{减少样本量并提高效能}{minimizing sample size and maximizing power}{82}
\idxentry{1}{检查对结局作用的研究}{studies of effect of testing on}{}
\idxentry{2}{分析}{analysis}{256}
\idxentry{2}{结局选择}{choice of outcome}{255}
\idxentry{2}{临床试验}{clinical trials}{256}
\idxentry{2}{决策分析}{decision analysis}{256}
\idxentry{2}{观察性研究}{observational studies}{255–256}
\idxentry{0}{结局变量}{Outcome variable}{5, 52, 418g}
\idxentry{1}{医学检查的准确度}{of medical tests}{246–247}
\idxentry{1}{数量}{number of}{199}
\idxentry{0}{提纲}{Outline}{}
\idxentry{1}{申请书提纲}{of proposal}{355}
\idxentry{1}{研究提纲}{of study}{13–14}
\idxentry{0}{过度诊断}{Overdiagnosis}{255, 419g}
\idxentry{0}{过拟合}{Overfitting}{243, 252, 419g}
\idxentry{0}{过度匹配}{Overmatching}{167, 419g}
\idxletter{P}
\idxentry{0}{P值操纵}{P-hacking}{175}
\idxentry{0}{$P$值}{P value}{58–59, 160, 419g}
\idxentry{0}{配对测量}{Paired measurements}{79–81, 419g}
\idxentry{0}{成对匹配}{Pairwise matching}{166, 419g}
\idxentry{0}{纸质问卷}{Paper-based questionnaires}{307}
\idxentry{0}{简约性}{Parsimony}{47}
\idxentry{0}{参与者}{Participant}{419g。另见研究参与者（Study participants）}
\idxentry{1}{基线变量}{baseline variables}{202}
\idxentry{1}{随机盲法试验}{randomized blinded trial}{}
\idxentry{2}{足够样本量与招募安排}{adequate sample size and planning recruitment}{201–202}
\idxentry{2}{临床试验选择标准}{selection criteria for clinical trials}{200–201, 201t}
\idxentry{1}{招募足够人数}{recruiting sufficient numbers of}{34–35}
\idxentry{0}{参与者偏倚}{Participant bias}{44。另见回忆偏倚（Recall bias）}
\idxentry{0}{参与者识别号}{Participant identification number}{340, 419g}
\idxentry{0}{参与者变异}{Participant variability}{42, 419g}
\idxentry{0}{患者人群}{Patient population}{27}
\idxentry{0}{儿科麻醉中的围术期虚拟现实}{Pediatric anesthesia, perioperative virtual reality for}{}
\idxentry{1}{背景与意义}{background and significance}{13}
\idxentry{1}{方法}{methods}{}
\idxentry{2}{数据分析}{data analysis}{13}
\sourcepage{770}
\idxentry{2}{测量}{measurements}{13}
\idxentry{2}{设计概述}{overview of design}{13}
\idxentry{2}{随机化与盲法}{randomization and blinding}{13}
\idxentry{2}{样本量估计}{sample size estimates}{13}
\idxentry{2}{研究参与者}{study participants}{13}
\idxentry{1}{具体目标}{specific aim}{13}
\idxentry{0}{同行评审}{Peer review}{362, 419g}
\idxentry{0}{符合方案分析}{Per-protocol analysis}{211–212, 213f, 419–420g}
\idxentry{0}{人时}{Person-time}{123, 131, 420g}
\idxentry{0}{I期试验}{Phase I trial}{206, 420g}
\idxentry{0}{II期试验}{Phase II trial}{206, 420g}
\idxentry{0}{III期（关键性）试验}{Phase III (pivotal) trial}{206, 420g}
\idxentry{0}{IV期试验}{Phase IV trial}{206, 420g}
\idxentry{0}{物理空间}{Physical space}{325–326}
\idxentry{0}{PI。参见主要研究者}{PI. See Principal investigator (PI)}{}
\idxentry{0}{先导研究}{Pilot study}{83, 196, 329, 420g}
\idxentry{1}{随机盲法试验}{randomized blinded trial}{206, 216}
\idxentry{0}{安慰剂对照}{Placebo control}{104–105, 420g}
\idxentry{0}{安慰剂效应}{Placebo effect}{196, 420g}
\idxentry{0}{剽窃}{Plagiarism}{106, 420g}
\idxentry{0}{平台试验}{Platform trial}{224, 420g}
\idxentry{0}{政治观点}{Political views}{189}
\idxentry{0}{多分类变量}{Polychotomous categorical variables}{40, 421g}
\idxentry{0}{癌基因中的多态性}{Polymorphisms in cancer genes}{189}
\idxentry{0}{人群}{Population}{4, 421g}
\idxentry{1}{临床人群与代表性人群}{clinical vs. representative}{30–31}
\idxentry{1}{研究}{research}{11}
\idxentry{1}{分层}{stratification}{168}
\idxentry{0}{基于人群的样本}{Population-based samples}{31, 421g}
\idxentry{2}{病例与对照样本}{of cases and controls}{142–143}
\idxentry{0}{阳性预测值}{Positive predictive value}{247}
\idxentry{0}{实证主义研究范式}{Positivist research paradigm}{265}
\idxentry{0}{资助后管理人员}{Post-award manager}{326, 421g}
\idxentry{0}{事后假设}{Post hoc hypothesis}{54, 60–62, 421g}
\idxentry{0}{后验分布}{Posterior distribution}{236, 421g}
\idxentry{0}{后验概率}{Posterior probability}{235, 250, 421g}
\idxentry{0}{随机化后因素}{Postrandomization factors}{198, 213}
\idxentry{0}{检验效能}{Power}{6, 56–57, 57t, 64, 421g}
\idxentry{1}{应用与例子}{applications and examples}{94}
\idxentry{1}{差异}{differences}{103–104}
\idxentry{1}{提高效能的策略}{strategies for maximizing}{}
\sourcepage{771}
\idxentry{2}{使用连续变量}{continuous variables, using}{78–79}
\idxentry{2}{使用结局}{outcome, using}{82}
\idxentry{2}{使用配对测量}{paired measurements, using}{79–81}
\idxentry{2}{使用精密变量}{precise variables, using}{81}
\idxentry{2}{使用不等组样本量}{unequal group sizes, using}{81–82}
\idxentry{0}{以实践为基础的研究网络}{Practice-based research networks}{283, 421g}
\idxentry{0}{资助前管理人员}{Pre-award manager}{326, 421g}
\idxentry{0}{干预前后设计。参见前后研究（试验）设计}{Pre-post design. See Before-after study (or trial) design}{}
\idxentry{0}{精密}{Precise}{39}
\idxentry{0}{精密度}{Precision}{9, 41–42, 44t, 160, 421–422g}
\idxentry{1}{评估}{assessing}{42}
\idxentry{1}{提高准确度的策略}{strategies for enhancing}{42–43, 43t}
\idxentry{1}{变量精密度与减少样本量、提高效能}{of variable, minimizing sample size and maximizing power}{81}
\idxentry{0}{临床前试验}{Preclinical trial}{206, 422g}
\idxentry{0}{预测效度}{Predictive validity}{46, 422g}
\idxentry{0}{预测值，阴性}{Predictive value, negative}{247, 422g}
\idxentry{0}{预测值，阳性}{Predictive value, positive}{247, 422g}
\idxentry{0}{预测因素}{Predictor}{52}
\idxentry{0}{预测变量}{Predictor variable}{5, 422g}
\idxentry{1}{医学检查的准确度}{of medical tests}{246}
\idxentry{0}{孕妇研究}{Pregnant women, research on}{103}
\idxentry{0}{初步数据}{Preliminary data}{359, 422g}
\idxentry{0}{预试}{Pretest}{314, 329, 422g}
\idxentry{0}{检验前概率。参见先验概率}{Pretest probability. See Prior probability}{}
\idxentry{0}{患病率}{Prevalence}{117, 117t, 422g}
\idxentry{1}{比}{ratio}{118}
\idxentry{0}{现用者}{Prevalent users}{129}
\idxentry{0}{既往收集标本与数据的研究}{Previously collected specimens and data, research on}{109}
\idxentry{0}{主要假设}{Primary hypothesis}{62, 422g}
\idxentry{0}{主键}{Primary key}{341, 341f, 422g}
\idxentry{0}{主要结局}{Primary outcome}{13, 62, 200, 422g}
\idxentry{0}{主要研究问题}{Primary research question}{22}
\idxentry{0}{主要研究者}{Principal investigator (PI)}{323, 354, 423g}
\idxentry{1}{主要研究者的诚信与领导力}{integrity and leadership of}{325}
\idxentry{0}{先验分布}{Prior distribution}{235, 237f, 423g}
\idxentry{0}{先验概率}{Prior probability}{61, 235, 250, 423g}
\idxentry{0}{囚犯研究}{Prisoners, research on}{103}
\idxentry{0}{私人基金会}{Private foundations}{364}
\idxentry{0}{私人信息}{Private information}{97}
\idxentry{0}{概率抽样}{Probability sampling}{31–32, 32f, 423g}
\idxentry{0}{职业利益冲突}{Professional conflicts of interest}{107, 423g}
\sourcepage{772}
\idxentry{0}{专业学会}{Professional societies}{364–365, 365f}
\idxentry{0}{项目官员}{Program officers}{353}
\idxentry{0}{项目可持续性}{Program sustainability}{283–284}
\idxentry{0}{项目概要或摘要}{Project summary or abstract}{356}
\idxentry{0}{倾向评分}{Propensity score}{172t, 173–175, 423g}
\idxentry{0}{申请书}{Proposal}{353, 423g}
\idxentry{1}{良好申请书的特征}{characteristics of good proposal}{361}
\idxentry{1}{组成要素}{elements of}{356t}
\idxentry{2}{行政部分}{administrative parts}{356–358}
\idxentry{2}{起始部分}{beginning}{355–356}
\idxentry{2}{最后组成部分}{final components of}{360–361}
\idxentry{2}{研究策略}{research strategy}{358–360, 359t}
\idxentry{2}{具体目标}{specific aims}{358}
\idxentry{1}{习题}{exercises for}{368}
\idxentry{1}{资助}{funding of}{361–362}
\idxentry{2}{基金会与专业学会资助}{foundations and professional societies, grants from}{364–365, 365f}
\idxentry{2}{机构内部支持}{intramural support}{366}
\idxentry{2}{NIH资助与合同}{NIH grants and contracts}{362–364, 363t, 364f}
\idxentry{2}{产业研究支持}{research support from industry}{365–366}
\idxentry{1}{撰写}{writing}{353–355}
\idxentry{0}{前瞻性队列研究}{Prospective cohort study}{4, 423g}
\idxentry{0}{前瞻性研究}{Prospective studies}{116}
\idxentry{1}{优缺点}{advantages and disadvantages}{121t}
\idxentry{1}{与回顾性研究的比较}{retrospective studies vs.}{119–122}
\idxentry{0}{受保护健康信息}{Protected health information}{299, 423g}
\idxentry{0}{保护因素}{Protective factor}{5, 423g}
\idxentry{0}{研究方案}{Protocol}{353, 424g}
\idxentry{1}{修订}{revisions}{}
\idxentry{2}{小幅修订}{minor}{329–330}
\idxentry{2}{实质修订}{substantive}{330}
\idxentry{0}{心理测量适当性}{Psychometric adequacy}{318, 424g}
\idxentry{0}{心理测量特征}{Psychometric characteristics}{317–318, 317t, 424g}
\idxentry{0}{心理测量等价性}{Psychometric equivalence}{318, 424g}
\idxentry{0}{发表偏倚}{Publication bias}{300, 424g}
\idxentry{0}{目的性抽样（选择）}{Purposeful sampling (or selection)}{265, 424g}
\idxletter{Q}
\idxentry{0}{QDA软件。参见定性资料分析软件}{QDA software. See Qualitative data analysis (QDA) software}{}
\idxentry{0}{定性方法}{Qualitative approaches}{}
\idxentry{1}{资料收集}{data collection}{274t}
\idxentry{2}{保密保护}{confidentiality protection}{272–273}
\idxentry{2}{研究现场与参与者}{field sites and study participants}{272}
\sourcepage{773}
\idxentry{2}{访谈提纲}{interview guides}{273}
\idxentry{2}{访谈／焦点小组}{interviews/focus groups}{272}
\idxentry{2}{样本量的确定}{sample size, determination of}{271–272}
\idxentry{1}{数据管理、分析与报告}{data management, analysis, and reporting}{}
\idxentry{2}{管理与分析过程}{process of managing and analyzing}{276–277}
\idxentry{2}{定性分析目标}{qualitative analysis, goals of}{274–275, 274t}
\idxentry{2}{定性资料分析软件}{qualitative data analysis software}{275–276}
\idxentry{2}{撰写与分享结果}{writing and sharing results}{277}
\idxentry{1}{习题}{exercises for}{279}
\idxentry{1}{目标}{goals of}{274–275, 274t}
\idxentry{1}{根据问题选择设计}{matching design to question}{270–271}
\idxentry{1}{定性方法}{qualitative methods}{}
\idxentry{2}{把握复杂概念}{capturing complex concepts}{269–270}
\idxentry{2}{探索新想法}{exploring new ideas}{270}
\idxentry{2}{开发新测量并改进既有测量}{new measurements development and improving existing ones}{269}
\idxentry{1}{定性研究}{qualitative research}{}
\idxentry{2}{研究方法}{approaches in}{268–268}
\idxentry{2}{方法比较}{comparing approaches}{268–269, 268t}
\idxentry{2}{定义}{defined}{263–267}
\idxentry{2}{与定量研究的比较}{quantitative vs.}{266–267t}
\idxentry{1}{与定量技术的比较}{quantitative techniques vs.}{269}
\idxentry{0}{定性资料分析（QDA）软件}{Qualitative data analysis (QDA) software}{275–276}
\idxentry{0}{定性研究}{Qualitative research}{}
\idxentry{1}{研究方法}{approaches in}{267–268}
\idxentry{1}{方法比较}{comparing approaches}{268–269, 268t}
\idxentry{1}{数据分析}{data analysis}{274t}
\idxentry{1}{研究中的资料收集}{data collection in}{274t}
\idxentry{1}{定义}{defined}{263–267}
\idxentry{1}{与定量研究的比较}{quantitative vs.}{266–267t}
\idxentry{0}{质量控制}{Quality control}{424g}
\idxentry{1}{临床程序}{clinical procedures}{332–333, 333t}
\idxentry{1}{多中心合作研究}{collaborative multicenter studies}{336}
\idxentry{1}{数据管理}{data management}{334–335, 335t}
\idxentry{1}{习题}{exercises for}{339}
\idxentry{1}{药物临床试验质量管理规范}{Good Clinical Practice}{331–332, 331t}
\idxentry{1}{实验室程序}{laboratory procedures}{333–334}
\idxentry{0}{质量协调员}{Quality coordinator}{332}
\idxentry{0}{质量改进}{Quality improvement}{297, 424g}
\idxentry{0}{Qualtrics}{}{346}
\idxentry{0}{定量研究}{Quantitative studies}{}
\idxentry{1}{把握复杂概念}{capturing complex concepts}{269–270}
\sourcepage{774}
\idxentry{1}{研究中的数据分析}{data analysis in}{274, 274t}
\idxentry{1}{研究中的资料收集}{data collection in}{273–274, 274t}
\idxentry{1}{与定量研究的比较}{quantitative vs.}{266–267t, 269}
\idxentry{0}{准实验（准随机）设计}{Quasi-experimental (quasi-randomized) design}{229, 424g}
\idxentry{0}{查询}{Query}{348, 424g}
\idxentry{0}{QuesGen}{}{346}
\idxentry{0}{根据问题选择设计}{Question, matching design to}{270–271}
\idxentry{0}{问卷}{Questionnaire}{7, 424g}
\idxentry{0}{自报问题与回答}{Questions and responses, self-reported}{}
\idxentry{1}{分支问题}{branching questions}{309–310}
\idxentry{1}{封闭式问题}{closed-ended question}{308}
\idxentry{2}{有序回答选项}{ordered response options for}{308–309}
\idxentry{2}{回答选项}{response options}{308}
\idxentry{1}{日记与日志}{diaries and logs}{313–314}
\idxentry{1}{多条目工具的格式设计}{multi-item instruments, formatting of}{312–313, 315}
\idxentry{1}{多条目量表}{multi-item scales}{311–312}
\idxentry{1}{开放式问题}{open-ended questions}{307–308}
\idxentry{1}{问题的最佳措辞}{optimal wording of questions}{310}
\idxentry{1}{问题设计的陷阱}{pitfalls in question design}{310–311}
\idxentry{1}{问题的时间范围}{question time frames}{311}
\idxletter{R}
\idxentry{0}{R类资助}{R awards}{362}
\idxentry{0}{雷洛昔芬心脏研究（RUTH）试验}{Raloxifene Use for The Heart (RUTH) trial}{330}
\idxentry{0}{随机误差}{Random error}{9, 43t, 56, 160, 424g}
\idxentry{0}{随机样本}{Random sample}{33, 424–425g}
\idxentry{0}{随机化}{Randomization}{6, 33, 196, 425g}
\idxentry{1}{匹配对的随机化}{of matched pairs}{204}
\idxentry{1}{随机盲法试验}{randomized blinded trial}{203–204}
\idxentry{0}{随机盲法试验}{Randomized blinded trial}{3, 4t}
\idxentry{1}{基线变量}{baseline variables}{}
\idxentry{2}{额外测量}{additional measurements}{202}
\idxentry{2}{保存的标本}{banked specimens}{202}
\idxentry{2}{参与者}{participants}{202}
\idxentry{1}{盲法}{blinding}{204–206}
\idxentry{1}{临床试验}{clinical trial}{}
\idxentry{2}{随访与方案依从性}{follow-up and adherence to the protocol}{207–209, 207t}
\idxentry{2}{意向治疗、符合方案和按实际治疗分析}{intention-to-treat, per-protocol, and as-treated analysis}{211–212, 212– 213f}
\idxentry{2}{监测}{monitoring}{209–211}
\idxentry{2}{传统临床机构之外}{outside of traditional clinical sites}{209}
\idxentry{2}{亚组分析}{subgroup analysis}{213–214}
\sourcepage{775}
\idxentry{1}{嵌入医疗服务的临床试验}{clinical trials embedded in health care}{206}
\idxentry{1}{例子}{examples of}{217–218}
\idxentry{2}{耐甲氧西林金黄色葡萄球菌菌血症的联合抗生素治疗}{combination antibiotic therapy for methicillin resistant Staphylococcus aureus bacteremia}{217}
\idxentry{2}{早期乳腺癌他莫昔芬后使用来曲唑}{letrozole after tamoxifen for early breast cancer}{217–218}
\idxentry{2}{胰腺切除伴或不伴常规腹膜内引流}{pancreas resection with/without routine intraperitoneal drainage}{218}
\idxentry{1}{习题}{exercises for}{219–220}
\idxentry{1}{操作手册目录}{operations manual table of contents for}{338}
\idxentry{1}{结局测量}{outcome measurements}{198–199}
\idxentry{2}{不良事件}{adverse events}{199–200}
\idxentry{2}{复合结局}{composite outcomes}{199}
\idxentry{2}{结局变量数量}{outcome variables, number of}{199}
\idxentry{1}{参与者选择}{participants selection}{}
\idxentry{2}{足够样本量与招募安排}{adequate sample size and planning recruitment}{201–202}
\idxentry{2}{临床试验选择标准}{selection criteria for clinical trials}{200–201, 201t}
\idxentry{1}{预试验}{pilot studies}{206, 216}
\idxentry{1}{随机化}{randomization}{203–204}
\idxentry{1}{选择干预与对照条件}{selecting intervention and control conditions}{197f}
\idxentry{2}{对照选择}{choice of control}{197–198}
\idxentry{2}{干预选择}{choice of intervention}{196–197}
\idxentry{2}{共同干预}{cointerventions}{198}
\idxentry{1}{新疗法监管审批试验}{trials for regulatory approval of new therapies}{206, 207t}
\idxentry{1}{利用既有资料收集机制}{using existing data collection mechanisms}{302–303}
\idxentry{0}{随机临床试验}{Randomized clinical trials}{108–109}
\idxentry{0}{随机质量改进试验}{Randomized quality improvement trial}{303, 425g}
\idxentry{0}{随机登记试验}{Randomized registry trial}{302, 425g}
\idxentry{0}{研究罕见结局的效率}{Rare outcomes, efficiency for}{140}
\idxentry{0}{率}{Rate}{123, 132t, 425g}
\idxentry{0}{率比}{Rate ratio}{132, 425g}
\idxentry{0}{合理怀疑}{Reasonable doubt}{55}
\idxentry{0}{回忆偏倚}{Recall bias}{44, 144, 425g}
\idxentry{0}{受试者工作特征曲线}{Receiver operating characteristic (ROC) curve}{76, 240, 248, 248f, 425g}
\idxentry{0}{记录}{Record}{340, 426g}
\idxentry{0}{招募}{Recruitment}{33, 426g}
\idxentry{1}{规划}{planning}{201–202}
\idxentry{1}{招募足够参与者}{recruiting sufficient numbers of participants}{34–35}
\idxentry{1}{代表性样本}{representative sample}{33–34}
\idxentry{1}{与保留参与者的关系}{and retention}{328–329}
\idxentry{0}{递归分割}{Recursive partitioning}{251, 426g}
\idxentry{0}{红肉摄入与结直肠癌}{Red meat consumption, and colorectal cancer}{159t, 189–190}
\idxentry{0}{REDCap。参见研究电子数据采集}{REDCap. See Research Electronic Data Capture (REDCap)}{}
\sourcepage{776}
\idxentry{0}{参照完整性}{Referential integrity}{342, 426g}
\idxentry{0}{登记系统}{Registry}{128, 294–295, 426g}
\idxentry{0}{登记来源对照}{Registry-based controls}{142}
\idxentry{0}{断点回归设计}{Regression discontinuity design}{169–170, 426g}
\idxentry{0}{均值回归}{Regression to the mean}{231, 426g}
\idxentry{0}{临床研究法规}{Regulations on clinical research}{95}
\idxentry{0}{法规事务经理}{Regulatory affairs manager}{332}
\idxentry{0}{新疗法的监管审批}{Regulatory approval, of new therapies}{206, 207t}
\idxentry{0}{关系数据库}{Relational database}{340, 426g}
\idxentry{0}{相对风险（RR）}{Relative risk (RR)}{70, 118, 132, 140}
\idxentry{0}{信度}{Reliability}{41, 312, 318。另见精密度（Precision）}
\idxentry{0}{远程试验}{Remote trials}{209}
\idxentry{0}{RePORT。参见研究项目在线报告工具}{RePORT. See Research Portfolio Online Reporting Tools (RePORT)}{}
\idxentry{0}{代表性样本}{Representative sample}{33–34, 426g}
\idxentry{0}{可重复性}{Reproducibility}{240, 243, 426g}
\idxentry{1}{分析}{analysis}{244–245}
\idxentry{1}{设计}{designs}{243–244}
\idxentry{0}{可重复性研究}{Reproducibility study}{243–245, 427g}
\idxentry{0}{研究}{Research}{}
\idxentry{1}{研究的解剖结构}{anatomy of}{3t}
\idxentry{2}{背景与意义}{background and significance}{3}
\idxentry{2}{设计}{design}{3–5, 4t}
\idxentry{2}{研究问题}{research question}{2–3}
\idxentry{2}{统计问题}{statistical issues}{6}
\idxentry{2}{研究参与者}{study participants}{5}
\idxentry{2}{变量}{variables}{5–6}
\idxentry{1}{儿童}{children}{102–103}
\idxentry{1}{临床影响}{clinical impact of}{111–112}
\idxentry{1}{定义}{defined}{96}
\idxentry{1}{豁免IRB审查的研究}{exempted from IRB review}{98t}
\idxentry{1}{快速审查}{expedited review}{98t}
\idxentry{1}{人体参与者。参见人体参与者研究}{human subjects. See Human subjects, research}{}
\idxentry{1}{研究兴趣}{interest}{17}
\idxentry{1}{使用去标识化剩余标本与数据}{with leftover de-identified specimens and data}{100}
\idxentry{1}{疟疾}{malaria}{111}
\idxentry{1}{不端行为}{misconduct}{105}
\idxentry{1}{新生儿血液标本}{neonatal blood specimens}{100–101}
\idxentry{1}{研究网络}{networks}{298}
\idxentry{1}{养老机构住民}{nursing home residents}{104}
\idxentry{1}{研究提纲}{outline of study}{13–14}
\idxentry{1}{向参与者付费}{participants, payment to}{111}
\sourcepage{777}
\idxentry{1}{难以理解研究风险与益处的人}{persons with difficulty understanding risks and benefits of research}{103}
\idxentry{1}{研究的生理机制}{physiology of}{6}
\idxentry{2}{因果推断}{causal inference}{8}
\idxentry{2}{设计研究}{designing study}{7, 7f, 11–12}
\idxentry{2}{研究错误}{errors of research}{8–10, 9f, 10t}
\idxentry{2}{实施研究}{implementing study}{8, 8f}
\idxentry{1}{孕妇}{pregnant women}{103}
\idxentry{1}{囚犯}{prisoners}{103}
\idxentry{1}{研究结构}{structure of}{6f}
\idxentry{1}{研究团队}{team}{323–325, 324t}
\idxentry{1}{转化研究}{translational}{}
\idxentry{2}{从临床到人群研究}{clinical to population research}{11}
\idxentry{2}{从实验室到临床研究}{laboratory to clinical research}{10–11}
\idxentry{0}{研究电子数据采集（REDCap）}{Research Electronic Data Capture (REDCap)}{346, 347}
\idxentry{0}{研究假设}{Research hypothesis}{6, 52, 427g}
\idxentry{0}{研究不端。参见科研不端}{Research misconduct. See Scientific misconduct}{}
\idxentry{0}{研究项目在线报告工具（RePORT）}{Research Portfolio Online Reporting Tools (RePORT)}{21}
\idxentry{0}{研究申请书。另见申请书}{Research proposal. See also Proposal}{}
\idxentry{1}{良好申请书的特征}{characteristics of good proposal}{361}
\idxentry{1}{组成要素}{elements of}{356t}
\idxentry{2}{行政部分}{administrative parts}{356–358}
\idxentry{2}{起始部分}{beginning}{355–356}
\idxentry{2}{最后组成部分}{final components of}{360–361}
\idxentry{2}{研究策略}{research strategy}{358–360, 359t}
\idxentry{2}{具体目标}{specific aims}{358}
\idxentry{1}{习题}{exercises for}{368}
\idxentry{1}{资助}{funding of}{361–362}
\idxentry{2}{基金会与专业学会资助}{foundations and professional societies, grants from}{364–365, 365f}
\idxentry{2}{机构内部支持}{intramural support}{366}
\idxentry{2}{NIH资助与合同}{NIH grants and contracts}{362–364, 363t, 364f}
\idxentry{2}{产业研究支持}{research support from industry}{365–366}
\idxentry{1}{撰写}{writing}{353–355}
\idxentry{0}{研究问题}{Research question}{2–3, 17, 427g}
\idxentry{1}{特征}{characteristics of}{}
\idxentry{2}{合乎伦理}{ethical}{21–22}
\idxentry{2}{可行}{feasible}{19–21}
\idxentry{2}{重要}{important}{21}
\idxentry{2}{新颖}{novel}{21}
\idxentry{1}{FINE标准}{“FINE” criteria for}{20t}
\idxentry{1}{导师的选择与合作}{mentor, choosing and working with}{22–23}
\idxentry{1}{来源}{origins of}{17–18}
\sourcepage{778}
\idxentry{2}{专业知识}{expertise}{18}
\idxentry{2}{医学或公共卫生}{medical or public health}{18}
\idxentry{2}{新想法与技术}{new ideas and techniques}{18–19}
\idxentry{1}{制定研究计划}{study plan, developing}{}
\idxentry{2}{主要与次要问题}{primary and secondary questions}{22}
\idxentry{0}{申请书的研究策略}{Research strategy, of proposal}{358–360, 359t}
\idxentry{0}{尊重个人}{Respect for persons}{95, 427g}
\idxentry{0}{受访者偏倚}{Respondent bias}{44}
\idxentry{0}{应答率}{Response rate}{34, 427g}
\idxentry{0}{保留参与者}{Retention}{328–329}
\idxentry{0}{回顾性病例对照研究}{Retrospective case-control studies}{138}
\idxentry{0}{回顾性队列研究}{Retrospective cohort study}{4, 427g}
\idxentry{0}{回顾性研究}{Retrospective studies}{116, 427g}
\idxentry{1}{与前瞻性研究的比较}{prospective studies vs.}{119–122}
\idxentry{0}{风险}{Risk}{131, 132t, 140, 155, 427g}
\idxentry{1}{尽量减少风险}{minimizing}{101–102}
\idxentry{1}{检查的风险}{of tests}{254}
\idxentry{0}{风险差}{Risk difference}{83, 132, 250, 427g}
\idxentry{0}{危险因素}{Risk factor}{5, 427g}
\idxentry{0}{风险模型}{Risk model}{427–428g}
\idxentry{0}{风险比}{Risk ratio}{70, 118, 132, 140, 250, 428g。另见相对风险（RR）（Relative risk (RR)）}
\idxentry{1}{与比值比的关系}{odds ratio and}{155}
\idxentry{0}{风险集}{Risk set}{146, 428g}
\idxentry{0}{ROC曲线。参见受试者工作特征曲线}{ROC curves. See Receiver operating characteristic (ROC) curves}{}
\idxentry{0}{RR。参见相对风险}{RR. See Relative risk (RR)}{}
\idxentry{0}{导入期}{Run-in period}{208–209, 428g}
\idxentry{0}{驱动变量}{Running variable}{168}
\idxentry{0}{RUTH试验。参见雷洛昔芬心脏研究试验}{RUTH trial. See Raloxifene Use for The Heart (RUTH) trial}{}
\idxletter{S}
\idxentry{0}{样本}{Sample}{26, 428g}
\idxentry{1}{定义}{defined}{27}
\idxentry{1}{基于人群的样本}{population-based}{31}
\idxentry{0}{研究参与者样本}{Sample of study participants}{5, 7}
\idxentry{0}{样本量}{Sample size}{6, 52, 64, 160, 428g}
\idxentry{1}{补充说明}{additional points}{}
\idxentry{2}{多个假设与事后假设}{multiple and post hoc hypothesis}{60–62}
\idxentry{2}{主要与次要假设}{primary and secondary hypothesis}{62}
\idxentry{2}{变异性}{variability}{59–60, 60f}
\idxentry{1}{应用与例子}{applications and examples}{94}
\idxentry{1}{确定}{determining}{201–202}
\idxentry{1}{等样本量组}{equal-sized group}{}
\sourcepage{779}
\idxentry{2}{用卡方统计量或Z检验比较二分类变量的比例}{chi-squared statistic or Z test to compare proportions of dichotomous variables}{87–89}
\idxentry{2}{用t检验比较连续变量均数}{t test to compare means of continuous variables}{86}
\idxentry{1}{应避免的错误}{errors to avoid}{83–84}
\idxentry{1}{习题}{exercises for}{64}
\idxentry{1}{固定样本量}{fixed}{77–78}
\idxentry{1}{假设。参见假设}{hypothesis. See Hypothesis}{}
\idxentry{1}{样本量不足}{inadequate}{257}
\idxentry{1}{信息不足}{insufficient information}{82–83}
\idxentry{1}{减少样本量并提高效能的策略}{minimizing sample size and maximizing power, strategies for}{}
\idxentry{2}{连续变量}{continuous variables}{78–79}
\idxentry{2}{使用结局}{outcome, using}{82}
\idxentry{2}{配对测量}{paired measurements}{79–81}
\idxentry{2}{使用精密变量}{precise variables, using}{81}
\idxentry{2}{使用不等组样本量}{unequal group sizes, using}{81–82}
\idxentry{1}{其他考虑与特殊问题}{other considerations and special issues}{}
\idxentry{2}{分类变量与计数变量}{categorical and count variables}{72–73}
\idxentry{2}{聚集样本}{clustered samples}{73}
\idxentry{2}{退出研究者}{dropouts}{72}
\idxentry{2}{等效性与非劣效试验}{equivalence and noninferiority trials}{74–75}
\idxentry{2}{匹配}{matching}{74}
\idxentry{2}{多变量调整与特殊统计分析}{multivariable adjustment and special statistical analysis}{74}
\idxentry{2}{生存分析}{survival analysis}{73}
\idxentry{1}{规划}{planning}{52}
\idxentry{1}{与效能的关系}{power, and}{94}
\idxentry{1}{定性研究设计}{qualitative study design}{271}
\idxentry{1}{分析性研究与实验的方法}{techniques for analytic studies and experiments}{65, 66t}
\idxentry{2}{卡方检验}{chi-squared test}{68–71, 87–89}
\idxentry{2}{相关系数}{correlation coefficient}{71–72, 90}
\idxentry{2}{t检验}{t test}{66–68, 86, 93}
\idxentry{1}{描述性研究的方法}{techniques for descriptive studies}{75–76}
\idxentry{2}{连续变量}{continuous variables}{76, 91}
\idxentry{2}{二分类变量}{dichotomous variables}{76–77, 92}
\idxentry{1}{基本统计原则}{underlying statistical principles}{55, 55t}
\idxentry{2}{\ensuremath{\alpha}、\ensuremath{\beta}与效能}{alpha, beta, and power}{56–57, 57t}
\idxentry{2}{备择假设的侧数}{alternative hypothesis, sides of}{58}
\idxentry{2}{效应量}{effect size}{56}
\idxentry{2}{P值}{P values}{58–59}
\idxentry{2}{统计检验}{statistical test}{59}
\idxentry{2}{I类与II类错误}{type I and type II errors}{56}
\idxentry{0}{抽样}{Sampling}{26, 428g}
\sourcepage{780}
\idxentry{1}{病例}{cases}{141–142, 141f}
\idxentry{1}{有关说明}{comments on}{33}
\idxentry{1}{对照}{controls}{142–143}
\idxentry{1}{作为抽样便利手段的匹配}{convenience, matching as}{167}
\idxentry{1}{与有向无环图的关系}{directed acyclic graphs and}{187}
\idxentry{1}{医学检查的准确度}{of medical tests}{245–246}
\idxentry{1}{非概率样本}{nonprobability samples}{33}
\idxentry{1}{概率样本}{probability samples}{31–32, 32f}
\idxentry{0}{抽样偏倚。参见选择偏倚}{Sampling bias. See Selection bias}{}
\idxentry{0}{抽样框}{Sampling frame}{31, 428g}
\idxentry{0}{SAS。参见统计分析系统}{SAS. See Statistical Analysis System (SAS)}{}
\idxentry{0}{量表}{Scale}{39, 40t, 428g}
\idxentry{1}{分类变量}{categorical variables}{40}
\idxentry{1}{选择}{choosing}{41}
\idxentry{1}{数值变量}{numerical variables}{40–41}
\idxentry{0}{学术素养}{Scholarship}{18}
\idxentry{0}{科研不端}{Scientific misconduct}{54, 95, 105–106, 428g}
\idxentry{0}{季节性变化}{Seasonal changes}{231}
\idxentry{0}{二次数据分析}{Secondary data analysis}{291, 429g}
\idxentry{0}{次要假设}{Secondary hypothesis}{62, 429g}
\idxentry{0}{次要研究问题}{Secondary research question}{22, 429g}
\idxentry{0}{数据管理安全}{Security, data management}{349}
\idxentry{0}{选择或入组标准}{Selection (or entry) criteria}{28}
\idxentry{1}{低剂量睾酮对安慰剂提高绝经女性性欲的临床试验}{clinical trial of low-dose testosterone vs. placebo to enhance libido in menopause}{29t}
\idxentry{1}{临床与代表性人群}{clinical vs. representative populations}{30–31}
\idxentry{1}{制定标准}{establishing}{29}
\idxentry{0}{选择偏倚}{Selection bias}{9, 131, 141–143, 141f, 161, 428g}
\idxentry{0}{选择标准}{Selection criteria}{200, 429g}
\idxentry{0}{自填式测量}{Self-administered measures}{319}
\idxentry{0}{自报测量}{Self-reported measure}{307, 429g}
\idxentry{1}{施测}{administering}{}
\idxentry{2}{实施平台}{administration platforms}{319}
\idxentry{2}{访谈方式}{interviewing approaches}{320}
\idxentry{2}{自填与访谈员施测}{self- vs. interviewer administration}{319}
\idxentry{1}{其他测量策略}{alternative measurement strategies}{320}
\idxentry{1}{用于多样人群}{for diverse populations}{318}
\idxentry{1}{习题}{exercises for}{322}
\idxentry{1}{使用既有测量}{existing measures, using}{314}
\idxentry{1}{设计新测量}{new measures, designing}{314}
\idxentry{1}{问题与回答。参见自报问题与回答}{questions and responses. See Questions and responses, self-reported}{}
\sourcepage{781}
\idxentry{1}{公开测量工具数据库示例}{sample databases of publicly available}{315–316t}
\idxentry{1}{选择与调整的步骤}{selecting and adapting, steps for}{}
\idxentry{2}{既有测量}{existing measures}{316}
\idxentry{2}{工具预试}{pretesting instruments}{318}
\idxentry{2}{既往施测方法与目标人群}{previous administration methods and target populations}{317}
\idxentry{2}{心理测量特征}{psychometric characteristics}{317–318, 317t}
\idxentry{2}{量表构建与评分的可解释性}{scale construction and score interpretability}{316–317}
\idxentry{2}{整套工具}{set of instruments}{318}
\idxentry{2}{所依据的概念}{underlying concepts}{316}
\idxentry{2}{变量或构念}{variables or constructs}{316}
\idxentry{1}{类型}{types of}{307}
\idxentry{0}{SEM。参见均数标准误}{SEM. See Standard error of the mean (SEM)}{}
\idxentry{0}{灵敏度}{Sensitivity}{76, 240, 247, 429g}
\idxentry{0}{敏感性分析}{Sensitivity analysis}{133, 429g}
\idxentry{0}{变化敏感性}{Sensitivity to change}{318, 429g}
\idxentry{0}{使用传感器的试验}{Sensors, trials using}{110–111}
\idxentry{0}{连续横断面调查}{Serial cross-sectional surveys}{119}
\idxentry{0}{共同结果}{Shared effect}{429g}
\idxentry{1}{条件化}{conditioning on}{175–176, 175f}
\idxentry{1}{有向无环图}{directed acyclic graphs}{184–185, 185f}
\idxentry{0}{意义}{Significance}{}
\idxentry{1}{研究策略中的意义部分}{section, of research strategy}{358}
\idxentry{1}{背景与意义}{background and}{3}
\idxentry{1}{统计学显著性}{statistical}{6}
\idxentry{0}{简单队列研究}{Simple cohort study}{124f}
\idxentry{1}{与较复杂队列研究的比较}{more complex vs.}{123t}
\idxentry{0}{简单假设}{Simple hypothesis}{53, 429–430g}
\idxentry{0}{简单随机样本}{Simple random sample}{31–32, 32f, 430g}
\idxentry{0}{问题的简明性}{Simplicity, of questions}{310}
\idxentry{0}{采用单一金标准}{Single gold standard, application of}{257–258}
\idxentry{0}{单患者试验}{Single patient trials}{226}
\idxentry{0}{怀疑态度}{Skeptical attitude}{18}
\idxentry{0}{怀疑性先验分布}{Skeptical prior distribution}{235, 430g}
\idxentry{0}{滚雪球抽样}{Snowball sampling}{268, 430g}
\idxentry{0}{社会与经济不利处境}{Social and economic disadvantages}{104–105}
\idxentry{0}{SOPs。参见标准操作规程}{SOPs. See Standard operating procedures (SOPs)}{}
\idxentry{0}{特殊统计分析}{Special statistical analysis}{74}
\idxentry{0}{具体目标}{Specific aims}{358, 430g}
\idxentry{0}{具体假设}{Specific hypothesis}{54}
\idxentry{0}{限定}{Specification}{165–166, 165t, 178, 430g}
\idxentry{0}{特异度}{Specificity}{76, 240, 247, 430g}
\sourcepage{782}
\idxentry{0}{疾病谱偏倚}{Spectrum bias}{241, 430g}
\idxentry{0}{疾病严重程度范围}{Spectrum of disease severity}{241–242}
\idxentry{0}{分样本验证}{Split sample validation}{252, 430g}
\idxentry{0}{SQL。参见结构化查询语言}{SQL. See Structured Query Language (SQL)}{}
\idxentry{0}{标准差}{Standard deviation}{40t, 42, 430g}
\idxentry{0}{均数标准误（SEM）}{Standard error of the mean (SEM)}{86, 430g}
\idxentry{0}{标准操作规程（SOPs）}{Standard operating procedures (SOPs)}{331, 430g}
\idxentry{0}{标准化}{Standardization}{336, 431g}
\idxentry{0}{标准化效应量}{Standardized effect size}{67–68, 431g}
\idxentry{0}{统计分析系统（SAS）}{Statistical Analysis System (SAS)}{346}
\idxentry{0}{统计问题}{Statistical issues}{6}
\idxentry{0}{统计程序}{Statistical program}{340, 431g}
\idxentry{0}{统计学显著性}{Statistical significance}{6, 431g}
\idxentry{0}{统计检验}{Statistical test}{59, 65, 66t, 431g}
\idxentry{0}{研究策略中的统计部分}{Statistics section, of research strategy}{359}
\idxentry{0}{指导委员会}{Steering committee}{336, 431g}
\idxentry{0}{主干与追问法}{Stem and probe approach}{273}
\idxentry{0}{阶梯楔形试验}{Stepped wedge trial}{228–229, 228f, 431g}
\idxentry{0}{保存的材料}{Stored materials}{47–48, 48t}
\idxentry{0}{分层}{Stratification}{171–173, 172t, 431g}
\idxentry{0}{分层区组随机化}{Stratified blocked randomization}{203, 431–432g}
\idxentry{0}{分层随机抽样}{Stratified random sampling}{32, 432g}
\idxentry{0}{关联强度}{Strength of association}{178}
\idxentry{0}{结构式访谈}{Structured interviews}{307}
\idxentry{0}{结构化查询语言（SQL）}{Structured Query Language (SQL)}{348}
\idxentry{0}{Student t检验。参见t检验}{Student’s t test. See t test}{}
\idxentry{0}{研究预算}{Study budget}{326}
\idxentry{0}{研究设计}{Study design}{7, 7f}
\idxentry{1}{操作手册}{operations manual}{11}
\idxentry{1}{研究提纲}{study outline}{11}
\idxentry{1}{研究方案}{study protocol}{11}
\idxentry{0}{研究实施}{Study implementation}{}
\idxentry{1}{覆盖分析}{coverage analysis}{327}
\idxentry{1}{资料收集平台}{data collection platforms}{327}
\idxentry{1}{数据库设计}{database design}{328}
\idxentry{1}{习题}{exercises for}{339}
\idxentry{1}{确定最终方案}{finalizing protocol}{}
\idxentry{2}{结束阶段}{closeout}{330–331}
\idxentry{2}{方案的小幅修订}{minor protocol revisions}{329–330}
\idxentry{2}{预试、正式演练与预试验}{pretests, dress rehearsals, and pilot studies}{329}
\idxentry{2}{方案的实质修订}{substantive protocol revisions}{330}
\sourcepage{783}
\idxentry{1}{伦理审查委员会批准}{institutional review board approval}{327}
\idxentry{1}{领导与团队建设}{leadership and team-building}{325}
\idxentry{1}{操作手册和表格设计}{operations manual and forms development}{327–328}
\idxentry{1}{招募与保留参与者}{recruitment and retention}{328–329}
\idxentry{1}{研究团队}{research team}{323–325, 324t}
\idxentry{1}{空间与设备}{space and equipment}{325–326}
\idxentry{1}{研究预算}{study budget}{326}
\idxentry{1}{研究启动阶段}{study start-up}{323}
\idxentry{0}{骨质疏松性骨折研究}{Study of Osteoporotic Fractures}{29, 83, 122}
\idxentry{0}{研究提纲}{Study outline}{13–14}
\idxentry{0}{研究参与者}{Study participants}{3, 5, 26, 26f}
\idxentry{1}{习题}{exercises for}{37–38}
\idxentry{1}{管理}{management}{274}
\idxentry{1}{招募}{recruitment}{}
\idxentry{2}{招募足够参与者}{recruiting sufficient numbers of participants}{34–35}
\idxentry{2}{代表性样本}{representative sample}{33–34}
\idxentry{1}{抽样}{sampling}{26f}
\idxentry{2}{抽样说明}{comments on sampling}{33}
\idxentry{2}{非概率样本}{nonprobability samples}{33}
\idxentry{2}{概率样本}{probability samples}{31–32, 32f}
\idxentry{1}{选择标准}{selection criteria}{28, 29t}
\idxentry{2}{临床与代表性人群}{clinical vs. representative populations}{30–31}
\idxentry{2}{制定标准}{establishing}{29}
\idxentry{1}{参与者一览表}{table of}{342–343f}
\idxentry{1}{术语与概念}{terms and concepts}{}
\idxentry{2}{人群}{populations}{27}
\idxentry{2}{获得研究参与者的方案}{protocol for acquiring study participants}{28, 28f}
\idxentry{2}{样本}{samples}{27}
\idxentry{2}{研究结果的推广}{study findings, generalizing}{27–28, 27f}
\idxentry{0}{研究计划}{Study plan}{22, 432g}
\idxentry{1}{制定计划与主要、次要问题}{developing, primary and secondary questions}{22}
\idxentry{0}{研究启动阶段}{Study start-up}{287, 323, 432g}
\idxentry{0}{亚组分析}{Subgroup analysis}{213–214, 432g}
\idxentry{0}{受试者。参见参与者}{Subject. See Participant}{}
\idxentry{0}{深入社区参与}{Substantial community engagement}{285}
\idxentry{0}{抑制}{Suppression}{176, 185, 432g}
\idxentry{0}{替代标志物（替代结局）}{Surrogate marker (or surrogate outcome)}{198, 432g}
\idxentry{0}{调查}{Survey}{20, 432g}
\idxentry{0}{SurveyMonkey}{}{346}
\idxentry{0}{生存分析}{Survival analysis}{73, 432–433g}
\idxentry{0}{对暴露的易感性}{Susceptibility to exposures}{168}
\sourcepage{784}
\idxentry{0}{系统误差}{Systematic error}{9, 161。另见偏倚（Bias）}
\idxentry{0}{系统综述}{Systematic review}{18, 31, 300, 433g}
\idxentry{0}{系统样本}{Systematic sample}{32, 433g}
\idxletter{T}
\idxentry{0}{$t$检验（Student $t$检验）}{t test (or Student’s t test)}{66–68, 433g}
\idxentry{1}{连续变量}{continuous variables}{86}
\idxentry{1}{配对测量}{with paired measurements}{80}
\idxentry{1}{使用与误用}{use and misuse of}{93}
\idxentry{0}{测量一览表}{Table of measurements}{343f}
\idxentry{0}{联合检测}{Tandem testing}{246}
\idxentry{0}{目标人群}{Target population}{7, 26, 27, 433g}
\idxentry{0}{目标试验}{Target trial}{129}
\idxentry{0}{教学}{Teaching}{19}
\idxentry{0}{团队建设}{Team-building}{325}
\idxentry{0}{时间变化}{Temporal changes}{231}
\idxentry{0}{时间标准}{Temporal criteria}{28}
\idxentry{0}{检测阴性对照组}{“Test-negative” control group}{143}
\idxentry{0}{按检验结果抽样}{Test-result-based sampling}{246, 433g}
\idxentry{0}{检查结果}{Test results}{5}
\idxentry{0}{主题编码}{Thematic coding}{433g}
\idxentry{0}{主题（资料）饱和}{Thematic (or data) saturation}{271, 433g}
\idxentry{0}{治疗性误解}{Therapeutic misconception}{99}
\idxentry{0}{问题的时间范围}{Time frames, question}{311}
\idxentry{0}{时间序列设计}{Time-series design}{80, 433g}
\idxentry{0}{事件发生时间资料}{“Time to event” data}{132}
\idxentry{0}{研究申请书题目}{Title, of research proposal}{355–356}
\idxentry{0}{取舍}{Trade-offs}{12}
\idxentry{0}{转化研究}{Translational research}{10, 433g}
\idxentry{1}{从临床到人群研究}{clinical to population research}{11}
\idxentry{1}{从实验室到临床研究}{laboratory to clinical research}{10–11}
\idxentry{0}{治疗}{Treatment}{5, 176。另见干预（Intervention）}
\idxentry{0}{三角互证}{Triangulation}{270, 434g}
\idxentry{0}{真阴性乳腺X线检查}{True-negative mammogram}{246}
\idxentry{0}{真阳性乳腺X线检查}{True-positive mammogram}{246}
\idxentry{0}{建立信任}{Trust building}{287}
\idxentry{0}{如实告知}{Truth-telling}{96, 434g}
\idxentry{0}{双样本$t$检验}{Two-sample t test}{80, 434g}
\idxentry{0}{双侧假设}{Two-sided hypothesis}{58, 434g}
\idxentry{0}{I类错误（假阳性）}{Type I error (false-positive)}{56, 160, 210, 434g}
\idxentry{0}{II类错误（假阴性）}{Type II error (false-negative)}{56, 434g}
\idxletter{U}
\sourcepage{785}
\idxentry{0}{非预期关联}{Unanticipated associations}{61}
\idxentry{0}{无偏}{Unbiased}{39}
\idxentry{0}{基本统计原则}{Underlying statistical principles}{52, 55, 55t}
\idxentry{1}{\ensuremath{\alpha}、\ensuremath{\beta}与效能}{alpha, beta, and power}{56–57, 57t}
\idxentry{1}{备择假设的侧数}{alternative hypothesis, sides of}{58}
\idxentry{1}{效应量}{effect size}{56}
\idxentry{1}{P值}{P values}{58–59}
\idxentry{1}{统计检验}{statistical test}{59}
\idxentry{1}{I类与II类错误}{type I and type II errors}{56}
\idxentry{0}{不当影响}{Undue influence}{99, 434g}
\idxentry{0}{不等组样本量}{Unequal group sizes}{}
\idxentry{1}{减少样本量并提高效能}{minimizing sample size and maximizing power}{81–82}
\idxentry{0}{无信息先验分布}{Uninformative prior distribution}{235, 434g}
\idxentry{0}{英国生物样本库}{United Kingdom Biobank}{295}
\idxentry{0}{美国人口普查局}{U.S. Census Bureau}{294}
\idxletter{V}
\idxentry{0}{含糊假设}{Vague hypothesis}{54}
\idxentry{0}{效度}{Validity}{39, 46–47, 318, 434g}
\idxentry{0}{变异性}{Variability}{59–60, 60f, 64, 65, 318, 434g}
\idxentry{1}{来源}{sources of}{242}
\idxentry{0}{变量}{Variables}{5–6}
\idxentry{1}{变量名称}{names}{343–344, 344f}
\idxentry{0}{方差}{Variance}{71, 434–435g}
\idxentry{0}{VAS。参见视觉模拟量表}{VAS. See Visual analog scale (VAS)}{}
\idxentry{0}{验证偏倚}{Verification bias}{257–258, 435g}
\idxentry{0}{虚拟试验}{Virtual trials}{209}
\idxentry{0}{视觉模拟量表（VAS）}{Visual analog scale (VAS)}{309, 435g}
\idxentry{0}{自愿同意}{Voluntary consent}{98–99}
\idxentry{0}{易受伤害性}{Vulnerability}{102}
\idxentry{0}{易受伤害的研究参与者}{Vulnerable research participants}{102–105, 435g}
\idxentry{1}{儿童}{children}{102–103}
\idxentry{1}{难以理解研究风险与益处的人}{persons with difficulty understanding risks and benefits of research}{103}
\idxentry{1}{权力差异}{power differences}{103–104}
\idxentry{1}{孕妇}{pregnant women}{103}
\idxentry{1}{囚犯}{prisoners}{103}
\idxentry{1}{社会与经济不利处境}{social and economic disadvantages}{104–105}
\idxletter{W}
\idxentry{0}{等待名单设计}{Wait-list design}{226, 435g}
\idxentry{0}{知情同意豁免}{Waiver of informed consent}{100, 100t}
\idxentry{0}{洗脱期}{Washout period}{225, 435g}
\idxentry{0}{两种饮食的减重}{Weight loss by two diets}{60f}
\sourcepage{786}
\idxentry{0}{组内设计}{Within-group design}{225, 435g}
\idxentry{0}{撰写申请书}{Writing proposals}{353–355}
\idxletter{Z}
\idxentry{0}{$Z$检验}{Z test}{68, 435g}
\idxentry{0}{Zoomerang}{}{346}
